\subsection{ACID}
\label{sec:acid}
\acrshort{acid} is the conventional notion of correctness (informal because it's an acronym):
\begin{itemize}
    \item \bi{Atomicity}: A group of operations must take place in their entirety or not at all
    \item \bi{Consistency}: Operations should take the database from a correct state to another correct state
    \item \bi{Isolation}: Concurrent execution of operations should yield predictable and correct results
    \item \bi{Durability}: The DB needs to remember the state it is in at all moments, even when failures occur
\end{itemize}

% TODO: Conflict Serializable histories

\inlinedefinition[Database state] is the actual values stored in a database at a given point in time (in memory and storage).

\inlinedefinition[Consistent state] is a database state that is the result of \textit{correctly} applying \textit{operations} to the DB.

\inlinedefinition[Correct application] The model of correct executions of operations is sequential executions.

\inlinedefinition[Operations] The operations for transactions are inserts, updates and deletion of tuples.

\inlinedefinition[Consistency] The assumption of correctness of transactions is based on two things:
\bi{(1)} the user doesn't apply wrong transactions and \bi{(2)} the DBMS has mechanisms to prevent incorrect modifications, such as violations of constraints, etc.

\inlinedefinition[Durability] ensures that the changes of a completed transaction are remembered by the system and can be recovered.
This is often more strict and DBMS will keep the DB st ate and a history of all transactions.

The transaction history is captured in the DB log. The state at a given moment is captured in snapshots.
The log and snapshots allow to go back and forward in the history of the database by undoing or redoing transactions from a given snapshot.
Thus, very similar in concept to what \texttt{git} uses.


\subsubsection{Atomicity}
\inlinedefinition[Atomicity] dictates that a transaction only takes effect if it can be executed in its entirety, or in other words,
a transaction has to be executed in its entirety or not at all.
The issue with this is that violations are costly to correct.

Thus, atomicity requires isolation. If a transaction can see the intermediate state created by another transaction, then its input is not guaranteed to be consistent,
implying that its output is not guaranteed to be consistent, either.

The DBMS get around this issue by isolating the transactions, such that they behave as if they were alone in the database.
This of course means that we need to know when a transaction is completed so we can \textit{commit} it, if there are no conflicts.
If there are conflicts, we have to resolve them, rerun the transaction (note that this can be very expensive) or to prevent them altogether.

The isolation is enforced through concurrency control mechanisms to prevent conflicts.
The canonical concurrency control mechanism in DBs is locking.
We use locking policies to ensure that the transactions are isolated.
